Conventions of Theorem \claimref{synthesis.rung-unit-is-frame-boost}: $v_{a} = \tfrac12(e_{a}^{-} + e_{a}^{+})$, $w_{a} = \tfrac12(e_{a}^{+} - e_{a}^{-})$, so $Q(v_{a}) = Q(w_{a}) = 0$, $v_{a}\cdot w_{b} = \tfrac12\delta_{ab}$ and $v_{a}w_{b} + w_{b}v_{a} = \delta_{ab}$ (Fact \claimref{standard-model.imported-tower-commuting-copies-fock-module}); the conjugation is $x^{*} = \theta(\bar{x})$ with $\theta$ fixing $e_{a}^{+}$ and negating $e_{a}^{-}$, so $\theta v_{a} = w_{a}$ and $v_{a}^{*} = w_{a}$. For $c = \alpha\cdot v + \beta\cdot w$ the complex-bilinear form is $Q(c) = 2\sum_{a}\alpha_{a}\beta_{a}\,(v_{a}\cdot w_{a}) = \alpha^{\mathsf T}\beta$, and $c^{*} = \bar{\beta}\cdot v + \bar{\alpha}\cdot w$, so $c^{*}\cdot c = \tfrac12(|\alpha|^{2} + |\beta|^{2})$. The slot generators $E_{ab} = v_{a}w_{b}$ act on vectors by commutator as the matrix unit $e_{a}e_{b}^{\mathsf T}$ on the $v$-coefficients and as $-e_{b}e_{a}^{\mathsf T}$ on the $w$-coefficients, so $X = \sum X_{ab}E_{ab}$ acts by $(\alpha, \beta) \mapsto (X\alpha, -X^{\mathsf T}\beta)$, which for $X \in \uu(2)$ is $(X\alpha, \bar{X}\beta)$; the group $U(2)$ acts by $(\alpha, \beta) \mapsto (X\alpha, \bar{X}\beta)$ and preserves $Q$ (since $X^{\mathsf T}\bar{X} = \overline{X^{\dagger}X} = \mathbf{1}$) and $\mathsf{N}$. The hypercharge $Y = \tfrac{I_{4}}{2}\mathbf{1}$ has weight $+\tfrac12$ on $v$ and $-\tfrac12$ on $w$, as in Proposition \claimref{standard-model.ten-transported-ambient-space-doublets-are}. The frozen point $I$ lies in the ground and the rung summand is orthogonal to it, so $Q(I + c) = \Lambda/3 + Q(c)$ and $I + c$ is null exactly when $c \in L$; by Theorem \claimref{synthesis.rung-unit-transverse-to-norm}(1) every lift $X_{\mu}(I)$ has rung component $c_{\mu} \in L$, and $c_{1}$ is the rung component of Proposition \claimref{standard-model.transported-frozen-point}(i).

(1) $X$ fixes $c$ iff $X\alpha = \alpha$ and $\bar{X}\beta = \beta$, that is $X\bar{\beta} = \bar{\beta}$. On $L$, $\alpha^{\mathsf T}\beta = -H^{2} \ne 0$, so $\alpha \ne 0 \ne \beta$. If $\alpha$ and $\bar{\beta}$ are independent they span $\C^{2}$ and $X = \mathbf{1}$. If $\bar{\beta} = \bar{\kappa}\alpha$, $X$ fixes the line of $\alpha$ pointwise and, being unitary, acts on its orthogonal line by a phase, which may be any phase: the stabiliser is a circle. Then $Q(c) = \alpha^{\mathsf T}\kappa\bar{\alpha} = \kappa|\alpha|^{2} = -H^{2}$, so $\kappa = -H^{2}/|\alpha|^{2}$ is real and negative. Put $u = \alpha/|\alpha|$ and choose $A \in U(2)$ with $Au = e_{1}$ (for instance the matrix with rows $\bar{u}^{\mathsf T}$ and $(-u_{2}, u_{1})$). Then $A\alpha = |\alpha|e_{1}$ and $\bar{A}\beta = \kappa\,\overline{A\alpha} = -(H^{2}/|\alpha|)e_{1}$, which is $c_{\lambda}$ with $\lambda = H^{2}/|\alpha| = |\beta|$. Conversely $c_{\lambda}$ has $\alpha = (H^{2}/\lambda)e_{1}$, $\beta = -\lambda e_{1} \in \C\bar{\alpha}$, and $U(2)$ preserves the alignment and the lengths. Since $|\beta|$ is $U(2)$-invariant, distinct $\lambda$ lie on distinct orbits. At $c_{\lambda}$ the stabiliser fixes $e_{1}$ and rotates $e_{2}$: it is generated by $I_{4}E_{55} = Y - T_{3}$, as $Y - T_{3} = \tfrac{I_{4}}{2}(\mathbf{1} - \sigma_{3}) = \operatorname{diag}(0, I_{4})$.

(2) With $Dc = dc + \sum_{A}g_{A}A^{A}T_{A}c$ the gauge-field quadratic forms are $g_{A}g_{B}\operatorname{Re}((T_{A}c)^{*}\cdot T_{B}c)$ and $g_{A}g_{B}Q(T_{A}c, T_{B}c)$, $Q(\cdot,\cdot)$ the polarisation. At $c_{\lambda} = a\,v_{4} + b\,w_{4}$ ($a = H^{2}/\lambda$, $b = -\lambda$ real), $T_{A}c = (aX_{A}e_{1}, b\bar{X}_{A}e_{1})$, so $\operatorname{Re}((T_{A}c)^{*}\cdot T_{B}c) = \tfrac12(a^{2} + b^{2})\operatorname{Re}(X_{A}^{\dagger}X_{B})_{11}$, and with $X_{A} = \tfrac{I_{4}}{2}\sigma_{A}$ ($\sigma_{0} = \mathbf{1}$ for $Y$), $X_{A}^{\dagger}X_{B} = \tfrac14\sigma_{A}\sigma_{B}$, whose $(1,1)$ entry has real part $\delta_{AB}$ for $A, B \in \{1,2,3\}$, $(\sigma_{3})_{11} = 1$ for the pair $(T_{3}, Y)$ and $1$ for $(Y, Y)$. This is $\tfrac14 \mathsf{N}(c_{\lambda})P_{0}$ with $P_{0} = \operatorname{diag}(1,1) \oplus \bigl(\begin{smallmatrix}1&1\\1&1\end{smallmatrix}\bigr)$, and the couplings make it $\tfrac14 \mathsf{N} P$. For the bilinear form, $Q(T_{A}c, T_{B}c) = ab\operatorname{Re}(X_{A}^{\mathsf T}\bar{X}_{B})_{11}$ and $X_{A}^{\mathsf T}\bar{X}_{B} = \overline{X_{A}^{\dagger}X_{B}}$ for anti-Hermitian $X_{A}$, giving $\tfrac14 ab P = \tfrac14 Q(c_{\lambda})P$. The block $\bigl(\begin{smallmatrix} g^{2} & gg' \\ gg' & g'^{2}\end{smallmatrix}\bigr)$ has eigenvalues $0$ and $g^{2} + g'^{2}$. Finally $\mathsf{N}(c_{\lambda}) = \tfrac12(H^{4}\lambda^{-2} + \lambda^{2}) = H^{2}\cosh 2\eta$, as in Theorem \claimref{synthesis.rung-unit-is-frame-boost}(4), and $Q(c_{\lambda}) = -\Lambda/3$.

(3) $T_{c}L = \{\delta : \alpha^{\mathsf T}\delta\beta + \delta\alpha^{\mathsf T}\beta = 0\}$ is a complex hyperplane, of real dimension $6$, and the orbit has dimension $4 - 1 = 3$ by (1). At $c_{\lambda}$ the condition reads $a\,\delta\beta_{4} + b\,\delta\alpha_{4} = 0$ with $\delta\alpha_{5}, \delta\beta_{5}$ free. In the plane $H_{4}$ the solutions are $z(a, -b)$, $z \in \C$; $z$ real gives $(a, -b) = -\lambda\,dc_{\lambda}/d\lambda$, and $z$ imaginary gives $I_{4}(a, -b) = 2T_{3}c = 2Yc$, a gauge direction. In $H_{5}$ the gauge directions are $T_{1}c = \tfrac{I_{4}}{2}(a, -b)$ and $T_{2}c = -\tfrac12(a, b)$ (coefficients of $(v_{5}, w_{5})$), so the physical directions there form a real plane. The real vector $dc_{\lambda}/d\lambda$ is orthogonal to the purely imaginary $T_{3}c$, and its $H_{4}$ support makes it orthogonal to $T_{1}c, T_{2}c$; it is fixed by $Y - T_{3} = \operatorname{diag}(0, I_{4})$, which acts only on the $H_{5}$ coefficients, by $I_{4}$ on $\delta\alpha_{5}$ and $-I_{4}$ on $\delta\beta_{5}$, so it acts on the physical plane in $H_{5}$ as a rotation with $(Y - T_{3})^{2} = -1$: charges $\pm1$. With $dc_{\lambda}/d\eta = \lambda\,dc_{\lambda}/d\lambda = -(H^{2}/\lambda)v_{4} - \lambda w_{4}$, $Q = (H^{2}/\lambda)\lambda = H^{2}$ and $\mathsf{N} = \tfrac12(H^{4}\lambda^{-2} + \lambda^{2}) = \mathsf{N}(c_{\lambda})$.

(4) $e^{tK}$ multiplies $\alpha$ by $e^{t}$ and $\beta$ by $e^{-t}$, so $e^{tK}c_{\lambda} = (e^{t}H^{2}/\lambda)v_{4} - e^{-t}\lambda w_{4} = c_{\lambda e^{-t}}$; the boost $B_{s}$ of Theorem \claimref{synthesis.rung-unit-transverse-to-norm} ($v \mapsto sv$, $w \mapsto w/s$ on $H_{4}$) is $e^{(\log s)K}$ there. $\mathsf{N}(e^{tK}c) = \tfrac12(e^{2t}|\alpha|^{2} + e^{-2t}|\beta|^{2})$ has derivative $e^{2t}|\alpha|^{2} - e^{-2t}|\beta|^{2} = 2m(e^{tK}c)$ and second derivative $2(e^{2t}|\alpha|^{2} + e^{-2t}|\beta|^{2}) > 0$, so it is strictly convex with its only minimum where $m = 0$. On the family $m(c_{\lambda}) = \tfrac12(H^{4}\lambda^{-2} - \lambda^{2}) = -H^{2}\sinh 2\eta$, strictly decreasing in $\lambda$, and zero exactly at $\lambda = H$, where $\mathsf{N} = H^{2}$.

(5) $\mathsf{N}^{2} - m^{2} = |\alpha|^{2}|\beta|^{2}$ identically, and $|Q|^{2} = |\alpha^{\mathsf T}\beta|^{2} = |\langle\bar{\alpha}, \beta\rangle|^{2} \le |\alpha|^{2}|\beta|^{2}$ by Cauchy-Schwarz, with equality iff $\beta \in \C\bar{\alpha}$. Corollary \claimref{standard-model.one-vev-vector-relations-breaking-pattern} gives, for $c = a\,v_{4} + b\,w_{4}$, tree entries $|a|$ on $(L, \nu^{c}), (Q, u^{c})$ and $|b|$ on $(L, e^{c}), (Q, d^{c})$; for $c_{\lambda}$, $|a| = H^{2}/\lambda$ and $|b| = \lambda$. Read through the conjugate instead, the same corollary applied to $c_{\lambda}^{*} = \bar{b}\,v_{4} + \bar{a}\,w_{4} = -\lambda v_{4} + (H^{2}/\lambda)w_{4}$ gives up-type entries $|\bar{b}| = \lambda$, while the down-type entries of $c_{\lambda}$ are $|b| = \lambda$; so the two are equal, with product $\lambda^{2} = |\beta|^{2} = \mathsf{N}(c_{\lambda}) - m(c_{\lambda})$, since $\mathsf{N} - m = |\beta|^{2}$ identically, and $\lambda^{2} = H^{2}e^{2\eta}$. The tractor norm, $ab = -H^{2}$, enters only the reading through $c$. Finally $c_{\lambda}^{*} = \theta(c_{\lambda}) = -\lambda v_{4} + (H^{2}/\lambda)w_{4} = -c_{H^{2}/\lambda}$; the conjugation fixes $Q$ on $L$ (it conjugates $Q$, which is real there), preserves $\mathsf{N}$, exchanges $|\alpha|$ and $|\beta|$ and so reverses $m$; and since $-\mathbf{1} \in U(2)$, the orbit of $-c_{H^{2}/\lambda}$ is that of $c_{H^{2}/\lambda}$.

The script builds $\Cl(2,2)$ in GA.py, reads $Q$ and the Hermitian product from GA.dot and the slot generators from GA.mul, and checks: the pair relations, the frame involution and the generator action; for three values of $\Lambda$ and five frames, $Q(c_{\lambda}) = -\Lambda/3$, a one-dimensional stabiliser spanned by $Y - T_{3}$, trivial stabilisers at two misaligned points of $L$, and, at six aligned points of $L$ with complex coefficients, an explicit unitary matrix carrying them to $c_{H^{2}/|\alpha|}$; the two gauge quadratic forms entry by entry for five frames, with the characteristic polynomial $t(t - g^{2})^{2}(t - g^{2} - g'^{2})$ of $P$ at $g = 1$, $g' = \tfrac12$ evaluated at five points; the normal space of dimension $3$, the neutral line, the charged plane with $(Y - T_{3})^{2} = -1$ on it, and $Q(dc/d\eta) = \Lambda/3$, $\mathsf{N}(dc/d\eta) = \mathsf{N}(c_{\lambda})$; the boost as $e^{(\log s)K}$ on the family, the closed forms of $\mathsf{N}$ and $m$ along it, the minimum at the balanced frame and $\mathsf{N}^{2} = Q^{2} + m^{2}$; and the product law read through $c$, the equal entries and the product $\mathsf{N} - m = \lambda^{2}$ read through $c^{*}$, the mirror $c_{\lambda}^{*} = -c_{H^{2}/\lambda}$, the identity $\mathsf{N}^{2} = |Q|^{2} + m^{2}$ at two aligned vectors off $L$ with complex $\kappa$ and the strict inequality at a misaligned vector.

Reading, argued and not part of the statement. (a) \emph{The quantum.} On $L$ the rung field has one neutral physical scalar, the rung unit $\eta$, and one charged scalar; the second neutral scalar and the pseudoscalar of a pair of doublets are removed by the complex constraint $Q(c) = -\Lambda/3$, and the tractor norm, which would be the radial mode of a doublet, is frozen. So the quantum of this field that plays the Higgs boson's part is a fluctuation of the frame boost, with Hermitian kinetic coefficient $\mathsf{N}(c_{\lambda})$, and for $\lambda \gg H$ the canonical field is approximately $\lambda$ itself. (b) \emph{The potential.} Any $U(2)$-invariant potential restricted to the charge-preserving part of $L$ is a function of $\eta$ alone by (1), and a conjugation-invariant one is even in $\eta$ by (5), stationary at the balanced frame. A hypercharge D-term $\tfrac12 g'^{2}(m - \xi)^{2}$ puts the minimum at $m = \xi$; without a constant $\xi$ it is the balanced frame, where $\mathsf{N} = H^{2}$ and the weak scale would be the Hubble rate. With $\xi$, the number $v/H$ is traded for $\xi/H^{2}$: nothing here fixes it, so the premise \claimref{synthesis.rung-unit-second-free-number} is neither derived nor refuted. (c) \emph{A tension.} With one coupling $y$ through the pairing (a single $10$) and the up-type masses read through $c$, the product law says every charge-preserving vacuum on $L$ has tree masses with $m_{u}m_{d} = y^{2}\Lambda/3$, of order $y^{2}(10^{-42}\,\mathrm{GeV})^{2}$, whatever the frame; the observed third family has $m_{t}m_{b} \approx 700\,\mathrm{GeV}^{2}$ (record values, from memory). So either the electroweak vacuum is not on $L$, and the Higgs's tractor norm is a second quantity of order $m_{t}m_{b}/y^{2}$ rather than the frozen point's $-\Lambda/3$, or the up-type masses are read through $c^{*}$, as the Standard Model reads them through the conjugate doublet, so that they too are Hermitian readings; in the second case $c^{*} = \bar{\beta}\cdot v + \bar{\alpha}\cdot w$ puts $|\beta| = \lambda$ on the up-type pairs and both masses scale with the frame.
